\documentclass[12pt]{article}
\usepackage{graphicx} % Required for inserting images
\usepackage{hyperref}
\usepackage{makecell}
\usepackage{makeidx}
\usepackage{eurosym}
\usepackage{amsmath}
\newcommand{\firstPage}{
    \thispagestyle{empty}
    \begin{figure}
    \centering
    \includegraphics[scale=0.7]{Swellfish_logo.png}
    \end{figure}
    \author{
        \date{}
        \href{mailto://swellfish14@gmail.com}{swellfish14@gmail.com} \\
    } 
}
\usepackage{hyperref}
\usepackage{array}
\usepackage{tabularx}
\usepackage{adjustbox}

\newcounter{verscount}
\setcounter{verscount}{0}
\newcommand{\addversione}[5]{
	\ifdefined\setversione
		\setversione{#1}
	\else\fi
	\stepcounter{verscount}
	\expandafter\newcommand%
		\csname ver\theverscount \endcsname{#1&#2&#3&#4&#5}
}

\newcommand{\listversioni}{
	\ifnum\value{verscount}>1
		\csname ver\theverscount \endcsname
		\addtocounter{verscount}{-1}
		\\\hline
		\listversioni
	\else
		\csname ver\theverscount \endcsname\\\hline
	\fi
}

\newcommand{\makeversioni}{
	\begin{center}
		\begin{tabularx}{\textwidth}{|c|c|X|X|X|}
		\hline
		\textbf{Versione} & \textbf{Data} & \textbf{Redattore} & \textbf{Verificatore} & \textbf{Descrizione} \\
		\hline
		\listversioni
		\end{tabularx}
	\end{center}
	\clearpage
}

\begin{document}

\graphicspath{ {../templates/img/} }

\title{Piano di progetto}

\firstPage

\maketitle

\begin{center}
    \begin{tabular}{r | l}
		\multicolumn{2}{c}{\textit{Informazioni}}\\
		\hline
		
			\textit{Redattori} &
			[Davide Porporati, Elena Marchioro, Francesco Naletto]\makecell{}\\

			\textit{Revisori} &
			[Jude Vensil Braceros]\makecell{}\\
			\textit{Responsabili} &
			[Andrea Veronese]\makecell{}\\
		      \textit{Uso} & 
                [Esterno]\makecell{}\\
    \end{tabular}
\end{center}

\begin{center}
    \textbf{Descrizione}\\
    File contenente il piano di qualifica. Contiene le metriche e i criteri di accettazione dei prodotti.
\end{center}

\pagebreak

\tableofcontents
\pagebreak

\printindex 

\addversione{0.0.0}{25/04/2023}{Andrea Veronese}{Davide Porporati}{Creata struttura di base del documento}
\addversione{0.0.1}{27/04/2023}{Davide Porporati, Elena Marchioro, Francesco Naletto}{Jude Vensil Braceros}{Creata struttura di base del documento}
\makeversioni

\section{Introduzione}
Lo scopo di questo documento è quello di fornire un inieme di strumenti atti a verificare se accettare o meno un prodotto, sia esso software o documentale.
\section{Analisi d'insieme}
\subsection{Qualità di prodotto}
\subsection{Documentazione}
 \begin{itemize}
	\item  i documenti devono essere scritti in modo chiaro e comprensibile per gli utenti.
	\item  i documenti devono essere scritti correttamente dal punto di vista grammaticale.
\end{itemize}
Per valutare la documentazione è stato concordato l'utilizzo dell'indice di Gulapease, che permette di valutare la leggibilità di un testo, tarando l'indice sulla lingua italiana.
\[ \text{Indice di Gulpease} = 89 + \frac{300*\text{\#frasi} - 10*\text{\#lettere}}{\text{\#parole}} \]


\begin{itemize}
	\item \#lettere: numero di caratteri alfanumerici;
	\item \#parole: numero di gruppi di caratteri alfanumerici;
	\item \#frasi: numero di gruppi di punti o punti e virgola consecutivi.
\end{itemize}
L'indice di Gulpease oscilla su una scala di valori tra 0 e 100, dove 100 indica la leggibilità più alta e 0 la leggibilità più bassa:\\
indice di Gulpease $<$ 80: difficoltà nel leggere per chi ha la licenza elementare \\
indice di Gulpease $<$ 60: difficoltà nel leggere per chi ha la licenza media
indice di Gulpease $<$ 40: difficoltà nel leggere per chi ha un diploma superiore
\subparagraph{Prodotti coinvolti:}
\begin{center}
	\begin{tabularx}{\textwidth}{|X|X|X|}
		\hline
		\textbf{Prodotto} & \textbf{Valore accettabile} & \textbf{Valore ottimale } \\
		\hline
		Documenti interni & $>$ 40                      & $>$ 50                    \\
		\hline
		Documenti esterni & $>$ 40                      & $>$ 60                    \\
		\hline
	\end{tabularx}\\[8pt]
	\mbox{}\\
\end{center}
\subsection{Software}
\subsubsection{Copertura statement}
Per valutare la copertuta degli statement verrà utilizzata una metrica basata sullo statement coverage. Con questa metrica tutti gli statement del codice sorgente verranno eseguiti almeno una volta.
\subparagraph{Prodotti coinvolti:}
\begin{center}
	\begin{tabularx}{\textwidth}{|X|X|X|}
		\hline
		\textbf{Prodotto} & \textbf{Valore accettabile } & \textbf{Valore ottimale } \\
		\hline
		Software          & $>$ 80\%                     & $>$ 95\%                     \\
		\hline
	\end{tabularx}\\[8pt]
	\mbox{}\\
\end{center}
\subsubsection{Copertura branch}
Un altra metrica che verrà utilizzata sarà quella del Branch coverage, con la quale si verificherà se tutti i branch del codice sono coperti da test.
Una copertura bassa andrà ad indicare che ci sono degli scenari dell'applicazione che mancano di test a sufficienza.

\subparagraph{Prodotti coinvolti:}
\begin{center}
	\begin{tabularx}{\textwidth}{|X|X|X|}
		\hline
		\textbf{Prodotto} & \textbf{Valore accettabile } & \textbf{Valore ottimale } \\
		\hline
		Software          & $>$ 80\%                     & $>$ 95\%                     \\
		\hline
	\end{tabularx}\\[8pt]
	\mbox{}\\
\end{center}
\end{document}
